\scanpage{263}
\mspage{4}

$I$, c'est donc le cardinal de l'ens \unsure{des|de} ses chaînes composables maximales.

Deux chaînes composables maximales sont toujours de \unsure{signe|signes} \unsure{opposé|opposés}, si l'une est directe l'autre est rétrograde, et inversement. \deleted{Si la pre} La première chaîne composable maximale \inserted{$C_1$} (disons simplement ``\uline{chaîne structurale}'') a le $\mathord{=}$ signe que $u_1$ ; \deleted{s'il est $+$} \deleted{i.e. si $u_1$ est directe,} on dit que la \uline{direction initiale} de $I$ est directe, resp. rétrograde, s'il en est ainsi de $u_1$, ou ce qui revient au $\mathord{=}$, de $C_1$.

De $\mathord{=}$ on définit la \uline{direction finale} de $I$ comme celle de $u_n$, ou encore celle de $C_r$. \deleted{Ce n'est que dans le} Direction initiale et finale sont égales si la longueur utile $r$ de $I$ est \inserted{im}paire, inégales si elle est paire. (\uline{\unsure{NB}} Ce n'est que pour $n=0$, i.e. quand $I=\{0\}$ est la cat. ponctuelle, que direction initiale et direction finale ne sont pas définies).

Les chaînes composables correspondent aux sous-catégories pleines de $I$ qui sont isomorphes à une catégorie $\mathbb{\Delta}_i$ (où $i$ est la longueur de la chaîne). Notons que $\mathbb{\Delta}_i$ n'a pas d'autre automorphisme que l'identité, donc les chaînes maximales \inserted{de longueur donnée $i\geq1$} correspondent exactement aux foncteurs

\diagram{\begin{tikzcd}[ampersand replacement=\&]
\mathbb{\Delta}_i\arrow[r,hook] \& I
\end{tikzcd}}

qui sont \deleted{\unsure{des mono}} injectifs sur les sommets. Un tel foncteur est \deleted{alors \unsure{automatiq}} pleinement fidèle, et un monomorphisme. \deleted{Une chaîne} Ce foncteur définit une application \inserted{\unsure{str|str.}} croissante ou \deleted{\unsure{détr.}} décroissante sur

% archive sheet number: 262
